\section{External Consistency Check: CADE Cartel Convictions}
\label{sec:cade}

Before turning to prices, we ask whether FL firms show up near known
cartels. CADE (Conselho Administrativo de Defesa Econ\^omica)
convicted 65 procurement cartel cases during 2009--2019; 47 convicted
firms match to BEC via CNPJ. Among the 2,735 FL firms, 193 (7.1\%)
co-participate with at least one CADE-convicted firm---3.5 times the
expected rate under a permutation null stratified by participation
quartile ($p < 0.001$; 1,000 iterations). That 92.9\% show no CADE
overlap is unsurprising given the low base rate of cartel prosecution
\citep{connor2007price}: convictions are the visible tip, and the
screen is designed to reach the rest.

\paragraph{Interpretation.}
Co-participation measures proximity to suspicious procurement
environments, not cartel membership. Once we condition on
participation count, the FL--CADE overlap rate becomes
indistinguishable from that of other always-losers ($p = 0.93$;
Section~\ref{sec:ground_truth}): the unconditional excess is driven
by high participation volume, which is the anomaly the screen is
designed to exploit. The AUC therefore reflects the screen's ability
to identify high-participation firms that operate near known cartel
environments, not its ability to distinguish colluders from
non-colluders conditional on participation intensity.
The permutation test (3.5$\times$, stratified
by participation quartile) indicates that the excess is not purely
mechanical, but within-quartile variation in participation counts
may still contribute to the association.
For enforcement purposes, this is the relevant target:
a triage rule that flags environments warranting closer scrutiny,
not a tool that adjudicates firm-level liability. The screen does
not identify which firms are colluding; it identifies where an
enforcement agency should look next.

\paragraph{Dual pattern.}
Most CADE-convicted firms in BEC are \emph{winners}---the designated
recipients of cartel rents. But three convicted firms are themselves
classified as frequent losers (Table~\ref{tab:cade_fl_firms}). All
three participated across multiple item groups and years without ever
winning---the defining FL pattern.

\begin{table}[htbp]
\centering
\caption{CADE-convicted firms classified as frequent losers}
\label{tab:cade_fl_firms}
\begin{adjustbox}{max width=\textwidth}
\footnotesize
\begin{tabular}{lllcc}
\toprule
Firm & Cartel & CADE process & Tenders & Wins \\
\midrule
Sol Tecnologia & Solar water heaters & 08012.001273/2010-24 & 84 & 0 \\
Nova Esperan\c{c}a Locadora & School transport.\ & 08700.005876/2019-85 & 97 & 0 \\
Jofran Com\'ercio & Trash bags (Op.\ Colludium) & 08700.005789/2015-02 & 65 & 0 \\
\bottomrule
\end{tabular}
\end{adjustbox}
\end{table}

\paragraph{Co-bidder analysis.}
FL firms appear in 5.6\% of CADE-firm tenders versus a 2.1\% baseline
(2.6$\times$). For individual convicted firms the rate is much
higher---up to 69\% for Mayfran (school transportation). Excluding
the three directly convicted FL firms and their 246 tender-items
leaves $\hat{\beta}$ unchanged at 0.064.

\paragraph{Beyond known cartels.}
Dropping all 31,447 CADE-involved tender-items barely moves the
coefficient: $\hat{\beta} = 0.062$ ($N = 1{,}622{,}954$) versus
0.064 in the full sample (Table~\ref{tab:excl_cade}). The price
association is not driven by tenders CADE has already flagged---which
is the screen's practical value: it reaches markets that existing
enforcement records do not.

